\section{System and Release Protocol}
\label{sec:system}
\label{sec:impl}
\subsection{A Release Crossing Two Availability Boundaries}
\label{sec:motivation}
Our running example has three packages: \code{assets} produces an npm build,
\code{base} packages that output into a Docker image, and \code{app} builds
from the versioned \code{base} image. Arrows in Fig.~\ref{fig:stages} run from
provider to consumer. The second builder obtains its base through a registry;
we do not assume that it shares the first builder's local image store.

\begin{figure*}[t]
\centering
\begin{tikzpicture}[>=Latex,node distance=16mm,
 every node/.style={font=\small},
 stage/.style={draw,rounded corners=2pt,minimum width=28mm,minimum height=9mm,align=center}]
\node[stage] (ab) at (0,0) {Build assets\\npm output};
\node[stage] (bb) at (6,0) {Build base\\Docker image};
\node[stage] (cb) at (12,0) {Build app\\\texttt{FROM base:v}};
\node[stage] (ap) at (0,-1.65) {Publish assets};
\node[stage] (bp) at (6,-1.65) {Publish base};
\node[stage] (cp) at (12,-1.65) {Publish app};
\draw[->] (ab) -- node[above,font=\scriptsize]{local output} (bb);
\draw[->] (ab) -- (ap);
\draw[->] (bb) -- (bp);
\draw[->,thick] (bp) -- node[above,sloped,font=\scriptsize]{registry required} (cb);
\draw[->] (cb) -- (cp);
\draw[->] (ap) -- node[below,font=\scriptsize]{publish order} (bp);
\draw[->] (bp) -- node[below,font=\scriptsize]{publish order} (cp);
\end{tikzpicture}
\caption{Stage dependencies, not a measured timeline. Publishing \code{assets}
may overlap building \code{base}. Building \code{app} must wait for
\code{base}'s publication. Availability modes are configured requirements:
an npm dependency fetched from a registry would also require publication.}
\label{fig:stages}
\end{figure*}

An npm label alone does not justify overlap: \code{base} must actually consume
local output. Similarly, Docker does not universally require a push between
builds; the registry barrier is needed for the isolated-builder consumption
mode in this example. The provider setting
\code{isBuildWaitingPublish: true} declares that its consumers require its
published artifact. Here it is set on \code{base}, while \code{assets} uses
the default local-build mode. Consumer publication is ordered after provider
publication in both modes. A missing stage script is a no-op stage, not a
reason to omit dependencies or completion records.

Suppose \code{assets} is at version 1.4.2, \code{base} at 0.8.2, and
\code{app} at 2.1.0. An intent to release an assets feature and patch all
transitive consumers plans versions 1.5.0, 0.8.3, and 2.1.1. After the first
two packages publish and are recorded, \code{app} fails. A planner that only
propagates from \emph{currently unreleased providers} can now lose
\code{app}: the original provider is no longer pending, so its intent reaches
nobody. The package is an orphan in the release ledger. Our planner instead
asks whether \code{app} has received that intent. Its own pending window is
unchanged, so a retry still owes version 2.1.1.

\subsection{From Repository Inputs to a Plan}
\label{sec:ccme}
The planner reads a fixed repository revision, package definitions, dependency
edges, commit messages, and release tags. Configuration supplies build and
publish commands and availability requirements; manifests can supply package
and dependency identities. The plan specifies selected packages, next
versions, and reasons before any external publication occurs. Changed
inputs require a new plan; a preview is not an authorization for arbitrary
later repository states.

Release intent extends Conventional Commits~\cite{cc}. For example,
\code{feat(assets)\^{}\^{}: add theme assets} directly requests a minor
release of \code{assets} and propagates patch intent to transitive consumers.
An unmarked unit affects its own scope only. A scope can explicitly name
packages or be derived from changed paths. A commit can contain multiple
independent units, each retaining its own source set and propagation depth.
The formal model consumes normalized units, not a particular parser syntax.
The complete language grammar is outside this paper~\cite{ccme2}.

Version bumps and prerelease channels are separate axes. A unit can request a
channel transition without requesting a bump, and channel resolution can
restrict which propagated bumps are meaningful. Cancellation and correction
operate on pending intent; they do not unpublish existing artifacts. An exact
version directive pins a release, while a hold retains work without admitting
the package to the current plan. These controls explain the predicates needed
by the model; their CLI forms are not additional scientific contributions.

\subsection{Execution, Recording, and Recovery}
For each selected package, the executor schedules version reconciliation,
build, and publish work. A consumer's version/build work waits for provider
builds, or for provider publications when their availability mode requires it.
Build and publish tasks have separate concurrency limits. Independent ready
tasks may proceed concurrently; deterministic planning does not require a
fixed temporal order for independent completion events.

Successful publication creates a per-package completion obligation. The model
records it immediately as a tag. The implementation also has configured
recording and finalization steps, whose failures must be reported; their
crash boundaries are not assumed atomic with an arbitrary publish command.
The next invocation reconstructs outstanding work from the resulting records.
This avoids a separate persistent recovery database, but depends on preserving
and making the release tags available. A lock prevents cooperating release
jobs from writing the same remote release concurrently. It does not make
registry operations transactional or prevent an administrator from rewriting
history outside the protocol.

Two failure cases must remain distinct. If a publish is known to have failed,
its obligation stays pending. If the external operation succeeds but its
response or completion record is lost, an absent tag cannot distinguish that
success from failure. The integration must reconcile the artifact's identity
or use a publisher that safely accepts repetition before recording completion.
Forward recovery does not by itself solve this uncertainty~\cite{helland12}.
The mathematical claims below refer to the recorded state and make their
retry assumptions explicit.

\subsection{Applicability and Practical Limits}
The graph may contain one package: the same record-and-retry discipline still
applies even though there is no propagation. A monorepo supplies several
packages under one history. In a polyrepo, a control repository records
submodule pointers and the release graph; its commits and tags provide the
single history and ledger used by the model. The linked repositories keep
their own histories. This is a construction over a control history, not an
atomic transaction across independent Git repositories.

The integration contract has four parts: stable package and version
identities; declared dependencies and availability requirements; commands
whose successful exit means their promised stage has completed; and durable,
trustworthy completion records. Repeated publication must either be safe or
reconciled when its outcome is uncertain. These requirements also admit a
specification released as versioned files: validation or rendering can be its
build stage, and publishing the bundle its publish stage. This is a common
protocol for conforming integrations, not an assertion that every existing
release workflow or version scheme fits the model. In particular, the model
below assumes an acyclic graph and SemVer-compatible versions.

\paragraph{Shared release logic across operating systems}
Commands permit different build ecosystems without embedding their build
engines in the planner. A game developer can work on Windows while Linux CI
builds a dedicated server; conversely, a Linux developer can use Windows CI
for native packaging. A shared configuration holds the graph, lifecycle, and
portable commands. Small Linux and Windows files reference it and select
\code{sh} or PowerShell, with separate referenced scripts for platform-specific
steps. For example, \code{cmake --build build} can be shared after each
platform configures its own build directory. This reuses orchestration logic;
it does not translate shell syntax or make build outputs portable.

Conditional and iteration commands can compose shared tasks and gates.
However, in 1.8.2 the conditional wrapper itself requires a POSIX shell;
a native Windows caller without it must select the platform configuration
directly. Toolchains and SDKs must exist on each execution host. These are
integration capabilities, not measured Windows/Linux equivalence or evidence
of game-development adoption. The experiments evaluate narrower properties.

\subsection{A Stable CI Entry Point}
The same package graph can drive named test tasks and the release protocol.
A generic CI job invokes those operations rather than enumerating packages
and duplicating their dependencies in the workflow. Adding a package within
an already provisioned environment changes package configuration and scripts;
it need not change the CI workflow. A new runner, credential, permission, or
system dependency may still require provisioning changes. Likewise, a
configured test gate must explicitly reject missing required tests and block
publication on failure: a shared graph alone is not a testing guarantee.

Affected-task execution already appears in systems such as Nx~\cite{nx}.
Here the practical benefit is using one package model for testing and
recoverable release execution across artifact kinds. Test and release selection may use different windows and policies;
sharing the graph avoids duplicated package lists without replacing CI policy.
